\documentclass[11pt]{article}
\usepackage[margin=1in]{geometry}
\usepackage{amsmath,amssymb,amsthm}
\usepackage{hyperref}
\usepackage{booktabs}
\usepackage{graphicx}
\hypersetup{hidelinks}

\newtheorem{theorem}{Theorem}
\newtheorem{lemma}[theorem]{Lemma}
\newtheorem{corollary}[theorem]{Corollary}
\newtheorem{proposition}[theorem]{Proposition}
\theoremstyle{definition}
\newtheorem{definition}{Definition}

\setlength{\parskip}{4pt plus 1pt minus 1pt}
\setlength{\parindent}{0pt}

\title{The Genome as Knowledge Graph:\\Governance, Scaling, and the Architecture of Multicellular Life}
\author{Patrick D.\ McCarthy}
\date{}

\begin{document}
\maketitle

\begin{abstract}
Papers~1--4 of this series derive the structural constraints any
persistent system must satisfy under bounded context: graph structure,
escalation, inseparability of knowledge and action, and encoding
permanence. Papers~5--8 validate these constraints empirically in
cancer genomics. This paper shows that the genome itself is such a
system. DNA is the serialized projection of a directed knowledge graph.
Epigenetic marks are the write-ahead log. Repair layers are the
governance architecture. We show that the number of distinct repair
mechanisms scales as $\log_{10}(n)$ where $n$ is organism cell count
($R^2 = 0.84$, $p = 1.8 \times 10^{-4}$), and that this relationship
is driven by scale, not by reproduction rate. The naked mole rat and
the mouse---same cell count, same layer count, 300-fold difference in
cancer rate---demonstrate that layer \emph{count} tracks architecture
while layer \emph{confidence} tracks outcome. Cancer is not uncontrolled
growth. It is a resource allocation failure in a bounded system: an
autoscaling policy that lost its controls, consuming fixed resources
that coordinated function requires. The genome invested seven layers
of error correction in its data and approximately one in its index.
CHIP is the invoice for that asymmetry.
\end{abstract}

% ------------------------------------------------------------------
\section{Introduction}
\label{sec:intro}
% ------------------------------------------------------------------

Every system in this series operates under the same constraint: bounded
active context. A knowledge graph in a software system has a finite
working set. A language model has a finite context window. The genome
has a finite cell, finite blood supply, finite immune capacity.

The previous papers derive what follows from this constraint and
validate those derivations in cancer data. This paper makes a stronger
claim: the genome is not merely subject to the same constraints. It
\emph{is} a knowledge graph under bounded context, and its architecture
follows the predictions exactly.

We proceed in four steps. Section~\ref{sec:projection} establishes
DNA as the serialized projection of a running knowledge graph.
Section~\ref{sec:governance} shows that repair layers scale
logarithmically with organism size. Section~\ref{sec:confidence}
separates architecture (layer count) from confidence (layer quality)
and identifies the measurement gap. Section~\ref{sec:autoscaling}
reframes cancer as an autoscaling failure in a bounded resource pool.

% ------------------------------------------------------------------
\section{DNA as Serialized Projection}
\label{sec:projection}
% ------------------------------------------------------------------

A cell is a running system. It maintains metabolic state, responds to
signals, executes repair, and coordinates with neighboring cells. The
cell's active knowledge---its runtime graph---is encoded in protein
concentrations, metabolite levels, chromatin accessibility, and
epigenetic marks. This is the Meta/Action/Knowledge triad from
Paper~1 \cite{McCarthy2026a}, instantiated in biochemistry.

DNA is not the system. DNA is the system's serialized reboot contract.
When the cell divides, the runtime is destroyed (denaturing dissolves
the chromatin structure that constitutes the working index). What
survives is the sequence---the data---and a partial copy of the
epigenetic state. The daughter cell must reconstruct the runtime from
this projection.

The three encoding layers from Paper~4 \cite{McCarthy2026d} map
directly:

\begin{center}
\begin{tabular}{lll}
\toprule
\textbf{Encoding layer} & \textbf{Paper~4} & \textbf{Genome} \\
\midrule
Most permanent & Axioms & DNA sequence \\
Intermediate & Validated propositions & Epigenetic marks \\
Most volatile & Working memory & Protein/metabolite state \\
\bottomrule
\end{tabular}
\end{center}

Different cell types access different subsets of the DNA sequence. A
hepatocyte decompresses liver-relevant genes; a neuron decompresses
brain-relevant genes. The sequence is identical. The index state
differs. Cell type \emph{is} index state. This is selective
decompression: the bounded context constraint means no cell can
maintain the full graph in working memory, so each cell maintains a
view.

Cancer mutations cluster in tissue-relevant genes because those are
the only genes decompressed in that tissue. An unexpressed gene cannot
be disrupted by a replication error in a context where it is never
replicated. The mutation spectrum per tissue type is a map of which
subgraph was active.

\subsection{The write-ahead log}

In database systems, the write-ahead log (WAL) records intended
changes before they are applied to the data files. If the system
crashes mid-write, the WAL enables recovery. Epigenetic marks serve
the same function.

\begin{center}
\begin{tabular}{lll}
\toprule
\textbf{Database concept} & \textbf{Biology} & \textbf{Survives replication?} \\
\midrule
RAM (query state) & Protein/metabolite levels & No \\
Write-ahead log & Epigenetic marks & Partially (DNMT1 copies) \\
Data files & DNA sequence & Yes (7 repair layers) \\
Backup/replica & Germline DNA & Yes + meiotic recombination \\
\bottomrule
\end{tabular}
\end{center}

DNMT1 follows the replication fork, reading methylation marks from
the parent strand and copying them to the daughter. This is WAL replay.
Histone modifications distribute roughly 50/50 to daughter strands,
with neighboring histones providing quorum-based reconstruction. This
is partial checkpoint recovery.

The WAL has approximately one layer of error correction (DNMT1
maintenance copying, partially backed by histone quorum). DNA has
seven. This asymmetry has consequences.

\subsection{CHIP as WAL buffer overflow}

Clonal hematopoiesis of indeterminate potential (CHIP) is the
age-related accumulation of somatic mutations in blood stem cells.
Approximately 80\% of CHIP mutations occur in three genes: DNMT3A,
TET2, and ASXL1 \cite{Jaiswal2014}. All three are methylation and
chromatin maintenance genes. They are the WAL replay machinery.

CHIP is not data corruption. The DNA sequence is intact. CHIP is
index corruption: the machinery that maintains epigenetic state after
replication is failing. Each replication event that proceeds with
degraded WAL replay loses more index state. The cell gradually
forgets its type specification.

In database terms: the WAL buffer is full. The system has two options.
Block new writes until the buffer drains (senescence---the cell stops
dividing). Or write without journaling (cancer---the cell divides
without maintaining epigenetic coherence). Fast, no stalling, but one
bad write and the table is corrupt.

The 80\% concentration in methylation genes tells you exactly which
service is the bottleneck. The genome invested seven layers of error
correction in its data and approximately one in its index. CHIP is
the invoice for that asymmetry.

% ------------------------------------------------------------------
\section{Governance Scales Logarithmically}
\label{sec:governance}
% ------------------------------------------------------------------

\subsection{The prediction}

Each repair layer provides multiplicative suppression of the base
mutation rate. Polymerase alone (without proofreading) produces errors
at approximately $10^{-6}$ per base per replication. Each subsequent
layer (proofreading, mismatch repair, base excision repair, nucleotide
excision repair, double-strand break repair, cell cycle checkpoints,
immune surveillance) catches a fraction of what the previous layer
missed.

A note on the taxonomy: the discretization into seven layers is a
modeling choice, not a biological given. Double-strand break repair
could be split into homologous recombination and non-homologous end
joining. Base excision repair could be subdivided by substrate. Some
authors count differently. What matters for the scaling argument is
not the exact count but that the number of \emph{qualitatively distinct}
repair mechanisms increases with organism complexity. Any reasonable
discretization (5--9 layers) yields the same logarithmic relationship
with cell count. The specific taxonomy here follows the conventional
textbook decomposition.

For an organism with $n$ cells to be viable, the total mutation burden
must remain bounded:

\[
\mu_{\text{base}} \cdot \frac{n}{10^k} \leq T
\]

where $\mu_{\text{base}}$ is the unsuppressed base rate, $k$ is the
number of governance layers, and $T$ is the maximum tolerable mutation
burden. Solving for $k$:

\[
k \geq \log_{10}\left(\frac{\mu_{\text{base}} \cdot n}{T}\right) \propto \log_{10}(n)
\]

The number of layers required grows logarithmically with cell count.

\subsection{The data}

We compiled repair layer counts and cell counts across ten organisms
spanning 15 orders of magnitude in cell count (Table~\ref{tab:organisms}).

\begin{table}[h]
\centering
\small
\begin{tabular}{lrrrl}
\toprule
\textbf{Organism} & \textbf{Cell count} & \textbf{Layers} & \textbf{Mut.\ rate} & \textbf{Layers present} \\
\midrule
\textit{E.\ coli} & $1$ & 2 & $5.4 \times 10^{-10}$ & Proof, MMR \\
\textit{S.\ cerevisiae} & $1$ & 3 & $3.3 \times 10^{-10}$ & +BER \\
\textit{C.\ elegans} & $959$ & 5 & $2.7 \times 10^{-9}$ & +NER, apoptosis \\
\textit{D.\ melanogaster} & $5 \times 10^4$ & 5 & $3.5 \times 10^{-9}$ & +NER, apoptosis \\
Zebrafish & $10^{10}$ & 6 & $7.5 \times 10^{-9}$ & +DSB, checkpoints \\
Naked mole rat & $10^{10}$ & 7 & $1.0 \times 10^{-9}$ & All 7, enhanced \\
Mouse & $10^{10}$ & 7 & $5.4 \times 10^{-9}$ & All 7 \\
Human & $3.7 \times 10^{13}$ & 7 & $5.0 \times 10^{-10}$ & All 7 \\
Elephant & $10^{15}$ & 7 & $3.0 \times 10^{-10}$ & All 7, 20$\times$ TP53 \\
Bowhead whale & $10^{15}$ & 7 & $2.0 \times 10^{-10}$ & All 7, enh.\ DSB \\
\bottomrule
\end{tabular}
\caption{Governance layers across organisms. Mutation rates are per
base pair per cell division. Layer count reflects distinct repair/checkpoint
mechanisms present, not quality or redundancy within each layer. Sources:
Drake 1991, Lynch 2010, Denver et al.\ 2004, Cagan et al.\ 2022.}
\label{tab:organisms}
\end{table}

Linear regression of governance layers on $\log_{10}(\text{cell count})$
yields:

\[
\text{layers} = 0.290 \times \log_{10}(n) + 3.245 \qquad
R^2 = 0.843, \quad p = 1.8 \times 10^{-4}
\]

The relationship is logarithmic, as predicted
(Figure~\ref{fig:governance}, Panel~A).

\begin{figure}[ht]
\centering
\includegraphics[width=\textwidth]{analysis/results/governance_scaling/governance_scaling.pdf}
\caption{Governance scaling across organisms.
\textbf{A.}~Repair layer count vs.\ $\log_{10}(\text{cell count})$
($R^2 = 0.843$, $p = 1.8 \times 10^{-4}$).
\textbf{B.}~Orders of magnitude of mutation rate suppression vs.\ cell
count. Suppression is roughly constant ($\sim$3 orders), indicating
that per-layer confidence, not layer count, explains the residual
variance.
\textbf{C.}~Mutation burden (rate $\times$ cell count) vs.\ cell count.
Red line: expected burden with no governance. Blue line: observed
trajectory with log-scaled governance. Without $\log(n)$ scaling,
elephants would carry $10^9$ per-division mutations; observed burden is
orders of magnitude lower.}
\label{fig:governance}
\end{figure}

\subsection{The confound: reproduction rate}

An obvious objection: organisms with more cells also reproduce more
slowly. Perhaps governance layers track reproduction rate, not cell
count. Elephants have fewer offspring than mice. The correlation might
be driven by life history rather than scale.

We tested four candidate predictors (Table~\ref{tab:confound},
Figure~\ref{fig:confound}).

\begin{table}[h]
\centering
\begin{tabular}{lrr}
\toprule
\textbf{Predictor} & $R^2$ & $p$ \\
\midrule
$\log_{10}(\text{cell count})$ & 0.843 & $1.8 \times 10^{-4}$ \\
$\log_{10}(\text{reproduction rate})$ & 0.628 & $6.3 \times 10^{-3}$ \\
$\log_{10}(\text{generation time})$ & 0.890 & $4.2 \times 10^{-5}$ \\
$\log_{10}(\text{lifespan})$ & 0.900 & $2.8 \times 10^{-5}$ \\
\bottomrule
\end{tabular}
\caption{Governance layer count regressed on candidate predictors.
Lifespan and generation time show higher $R^2$ but are confounded
with cell count. The naked mole rat breaks the confound.}
\label{tab:confound}
\end{table}

Lifespan ($R^2 = 0.900$) and generation time ($R^2 = 0.890$) outperform
cell count ($R^2 = 0.843$) as raw correlates. With $N = 10$ and only
one overlapping pair (mouse/naked mole rat), this difference is not
statistically decisive. The honest statement: we cannot statistically
reject lifespan as the primary predictor from these data alone.

The argument for cell count as the architectural driver is not purely
statistical. It is mechanistic. The naked mole rat exposes the logic.

\begin{figure}[ht]
\centering
\includegraphics[width=\textwidth]{analysis/results/governance_scaling/confound_analysis.pdf}
\caption{Confound analysis: governance layers regressed on four
candidate predictors. Red circles highlight the mouse and naked mole
rat---same cell count, same layer count, different reproduction rate.
In the cell count panel (upper left), they overlap. In every other
panel, they separate. Cell count predicts layer \emph{count}.
Lifespan and generation time conflate count with \emph{confidence}.}
\label{fig:confound}
\end{figure}

\subsubsection{The naked mole rat}

The naked mole rat (\textit{Heterocephalus glaber}) and the laboratory
mouse have approximately the same cell count ($\sim 10^{10}$) and the
same seven governance layers. They differ in everything else:

\begin{center}
\begin{tabular}{lrr}
\toprule
\textbf{Metric} & \textbf{Mouse} & \textbf{Naked mole rat} \\
\midrule
Cell count & $10^{10}$ & $10^{10}$ \\
Governance layers & 7 & 7 \\
Mutation rate (per bp/div) & $5.4 \times 10^{-9}$ & $1.0 \times 10^{-9}$ \\
Generation time (days) & 60 & 365 \\
Lifespan (years) & 2 & 30 \\
Cancer rate & 30\% & $<0.1\%$ \\
\bottomrule
\end{tabular}
\end{center}

Same architecture. Same layer count. 300-fold difference in cancer
rate. Layer count tracks cell count (both $10^{10}$, both 7 layers).
Cancer rate tracks something else. That something is layer confidence.

% ------------------------------------------------------------------
\section{Architecture and Confidence}
\label{sec:confidence}
% ------------------------------------------------------------------

\subsection{The two-dimensional model}

Governance is not a scalar. It has two dimensions:

\begin{enumerate}
\item \textbf{Architecture}: the number of distinct governance layers.
   This is the structural decision. It scales as $\log(n)$ where $n$
   is cell count. It is set at speciation.
\item \textbf{Confidence}: the reliability of each layer. This is
   the investment decision. It varies within an architecture. It is
   set by selection pressure on lifespan.
\end{enumerate}

Total governance capacity is the product of these dimensions:

\[
G = \prod_{i=1}^{k} (1 - \epsilon_i)
\]

where $k$ is the layer count and $\epsilon_i$ is the error rate of
layer $i$. This product is dominated by the weakest layer.

\subsection{The $n$-nines problem}

In distributed systems, overall availability is bounded by the least
reliable component. A service mesh with six components at 99.99\%
availability and one component at 99\% delivers 99\% system
availability. The weak link is the bottleneck.

The same constraint governs biological governance. Seven repair layers
at $1 - \epsilon$ each yield system confidence
$(1-\epsilon)^7 \approx 1 - 7\epsilon$ for small $\epsilon$. But if
one layer is at $1 - 10\epsilon$ while the others are at
$1 - \epsilon$, the system confidence is approximately
$1 - 16\epsilon$, dominated by the weak layer.

Cancer does not need to defeat the governance system. It needs to
defeat \emph{one layer}. The cascade handles the rest.

To make this concrete: suppose each layer catches fraction $d_i$ of
errors that reach it. Layer 1 (proofreading) sees the raw error rate
$\mu_{\text{base}}$ and passes through $\mu_{\text{base}}(1 - d_1)$.
Layer 2 (MMR) sees that residual and passes through
$\mu_{\text{base}}(1-d_1)(1-d_2)$. The observed mutation rate is:

\[
\mu_{\text{obs}} = \mu_{\text{base}} \prod_{i=1}^{k} (1 - d_i)
\]

For humans, $\mu_{\text{base}} \approx 10^{-6}$ and
$\mu_{\text{obs}} \approx 5 \times 10^{-10}$, so:

\[
\prod_{i=1}^{7} (1-d_i) = \frac{5 \times 10^{-10}}{10^{-6}}
= 5 \times 10^{-4}
\]

If all layers had equal detection rate $d$, then $(1-d)^7 = 5 \times
10^{-4}$, giving $d \approx 0.65$ (each layer catches 65\% of errors
reaching it). But the product is dominated by the minimum. If six
layers are at $d = 0.90$ and one is at $d = 0.30$:

\[
(1 - 0.90)^6 \times (1 - 0.30) = 10^{-6} \times 0.70
= 7 \times 10^{-7}
\]

That single weak layer degrades the system by three orders of
magnitude. This is the $n$-nines problem: you cannot achieve five
nines of system reliability from a component running at three nines.
The system is bounded by its weakest layer.

\subsection{Measuring confidence per layer}

The decomposition above requires per-layer detection rates $d_i$.
These exist in fragments:

\begin{center}
\small
\begin{tabular}{llr}
\toprule
\textbf{Layer} & \textbf{Measurement method} & \textbf{Est.\ $d_i$} \\
\midrule
Polymerase proofreading & In vitro fidelity assays & $\sim$0.99 \\
Mismatch repair (MMR) & MMR$^{-/-}$ vs.\ wild-type rate & $\sim$0.99 \\
Base excision repair & 8-oxoG accumulation & $\sim$0.99 \\
Nucleotide excision repair & UV dose-response & $\sim$0.90 \\
DSB repair (HR/NHEJ) & HR: high fidelity; NHEJ: lossy & $\sim$0.85 \\
Checkpoints (p53/Rb) & Apoptotic rate per damage signal & $\sim$0.70 \\
Immune surveillance & NK/T-cell clearance of transformed cells & Unknown \\
\bottomrule
\end{tabular}
\end{center}

A caveat on the cascade model: the layers are not purely serial. BER
handles oxidative damage, NER handles bulky adducts, DSB repair handles
strand breaks. They are partially parallel, each addressing a different
class of lesion. The serial product formula overstates total suppression
because it treats each layer as if it operates on the same error stream.
The true architecture is a mix of serial (proofreading $\to$ MMR on
replication errors) and parallel (BER, NER, DSB on damage-type-specific
lesions), with checkpoints and immune surveillance as system-level
catches operating on all error classes.

The measurable quantity is the system output: proofreading and MMR
together suppress replication errors by $\sim 10^{4}$-fold (from
$10^{-6}$ to $\sim 10^{-10}$). The remaining layers primarily
address non-replication damage and transformed-cell clearance. Their
contribution to the per-base-per-division rate is smaller; their
contribution to preventing cancer is not.

The mouse/naked mole rat comparison illustrates the point without
requiring exact per-layer decomposition. Both have seven layers. Both
have the same replication fidelity machinery. The naked mole rat's
advantage is in the system-level layers: enhanced contact inhibition
(growth arrest at lower cell densities), more stable epigenetic
maintenance, and more robust p53 response. The 5.4-fold lower
mutation rate ($1.0 \times 10^{-9}$ vs.\ $5.4 \times 10^{-9}$)
is modest. The 300-fold cancer rate difference is not. The gap between
mutation rate difference and cancer rate difference is exactly the
signal from checkpoint and surveillance layers---the layers the
per-base rate does not capture.

The honest statement: we can measure architecture (layer count). We
can measure system output (mutation rate). We can partially decompose
via knockout experiments (MMR$^{-/-}$ mice) and cancer spectrum
analysis (MSI-high $=$ MMR failure, BRCA-mutant $=$ DSB failure,
TP53-mutant $=$ checkpoint failure). Each sequenced tumor is a
natural knockout experiment---a postmortem report identifying which
layer went down. But a full per-layer confidence matrix across species
does not yet exist.

\subsection{Where the data goes}

The naked mole rat's advantage is not in layer count. It is in
checkpoint confidence. High-molecular-weight hyaluronan triggers
early contact inhibition---growth arrest at lower cell densities
than in mice \cite{Seluanov2009}. This is the same checkpoint
layer (growth control), with higher detection confidence.

The elephant's advantage is not in layer count. It is in checkpoint
redundancy. Twenty copies of TP53 versus human two \cite{Sulak2016}.
Same layer. More copies. Higher confidence at the bottleneck.

Both organisms invested in the \emph{same} layer: the one most
likely to be the system's weak link. The investment is not random.
It targets the minimum.

% ------------------------------------------------------------------
\section{Cancer as Autoscaling Failure}
\label{sec:autoscaling}
% ------------------------------------------------------------------

\subsection{The base rate}

A bacterial cell mutates at approximately $10^{-6}$ to $10^{-3}$ per
base per generation. This is the base rate. It was not selected for.
It is the intrinsic error rate of polymerase-based replication.

Every human cell carries the same polymerase machinery. The base rate
is present in every one of the $3.7 \times 10^{13}$ cells. The seven
governance layers suppress it to $\sim 5 \times 10^{-10}$: a
$\sim$2000-fold reduction. The mutation rate is not something cancer
cells ``speed up to.'' It is the base rate that governance suppresses.
Cancer cells shed governance layers, and the base rate re-emerges.

Cagan et al.\ showed that somatic mutation rates scale inversely with
lifespan across 16 mammalian species, converging on similar end-of-life
mutation burdens \cite{Cagan2022}. This is governance scaling in action:
longer-lived organisms invest more in per-layer confidence, keeping the
lifetime accumulation bounded despite more cell divisions.

MSI-high cancers (mismatch repair lost): $\sim 10^{-7}$. Hypermutated
cancers (multiple layers lost): $\sim 10^{-5}$. The mutation rate is
a readout of remaining governance layers. It is the system's monitoring
dashboard: you can read how many services are still running from the
error rate.

\subsection{The autoscaling analogy}

Bacteria solved the scaling problem horizontally: stateless, fungible,
$O(1)$ governance per instance (Paper~1 \cite{McCarthy2026a}). No
$\log(n)$ governance overhead, which is why they won the biomass war.
But horizontal scaling cannot build an eye.

Multicellular life vertically scaled, and every failure mode of
distributed systems followed: consensus failures (defective signaling),
state corruption (epigenetic drift), resource contention (cachexia).

Every cell in a human body is horizontally scaling. Gut lining
replaces every 5 days. Skin every 2 weeks. Blood every 120 days.
The autoscaler runs constantly. This is the mechanism, not the
failure.

Cancer is not a cell that learned to replicate. It is a cell whose
autoscaling policy broke. The scale-up signal fires. The scale-down
signal does not.

\begin{center}
\small
\begin{tabular}{lll}
\toprule
\textbf{Autoscaling concept} & \textbf{Biology} & \textbf{Cancer} \\
\midrule
Scale-up trigger & Growth factors (EGF) & Constitutive activation \\
Scale-down trigger & Contact inhibition, TGF-$\beta$ & Lost \\
Max instances & Hayflick limit (telomeres) & Bypassed (telomerase) \\
Health check & p53 checkpoint & Disabled (TP53 mutated) \\
Circuit breaker & Apoptosis & Evaded (BCL2 overexpressed) \\
Resource quota & Nutrient dependence & Self-sufficient (Warburg) \\
Graceful shutdown & Anoikis & Ignored (metastasis) \\
\bottomrule
\end{tabular}
\end{center}

Every hallmark of cancer \cite{Hanahan2011} is a broken autoscaling
control. Not one is ``cell learned to replicate.'' The capacity was
always there. That is the base rate. The controls made it a managed
service instead of a runaway process.

\subsection{Bounded context makes misallocation fatal}

In an unbounded system, a misconfigured autoscaler is a nuisance. In
a bounded system, it is fatal. The organism has a fixed context
window: fixed blood supply, fixed glucose, fixed immune capacity.

Every resource consumed by uncoordinated growth is unavailable to
coordinated function. A tumor does not kill by being a tumor. It kills
by consuming glucose (Warburg effect), co-opting blood supply
(angiogenesis), occupying physical space, and hijacking immune
resources. This is noisy neighbor syndrome in a fixed resource pool.

The Warburg effect illustrates the point \cite{VanderHeiden2009}.
Cancer cells switch to glycolysis: 2 ATP per glucose instead of 36
from oxidative phosphorylation. Why choose the less efficient pathway? Because
glycolysis is \emph{faster}. In a bounded resource environment, the
service that captures resources fastest wins, even if it wastes 94\%
of them. A misconfigured autoscaler spinning up 100 instances at 2\%
utilization starves the service next to it that needs 3 instances at
90\%.

Cachexia---the wasting syndrome that kills most cancer patients---is
the system running out of resources. Not because the tumor consumed
them all directly, but because the immune response to the tumor
burned the budget. The system crashed not from the attack but from
the incident response. That is a DDoS from inside the mesh.

The tumor does not need to be malicious. It only needs to be
uncoordinated. In a bounded system, uncoordinated scaling
\emph{is} an attack on everything else.

% ------------------------------------------------------------------
\section{Governance Confidence Predicts Lifespan}
\label{sec:cagan}
% ------------------------------------------------------------------

Cagan et al.\ measured somatic mutation rates (single-base
substitutions per year in colonic crypts) across 16 mammalian species
\cite{Cagan2022}. The relationship between mutation rate and lifespan
is almost exactly the inverse predicted by the governance model
(Figure~\ref{fig:cagan}).

\begin{figure}[ht]
\centering
\includegraphics[width=\textwidth]{analysis/results/governance_scaling/cagan_kill_shot.pdf}
\caption{The Cagan relationship and its interpretation.
\textbf{A.}~Somatic mutation rate vs.\ lifespan across 16 mammalian
species (log-log). Slope $= -0.81$ ($R^2 = 0.84$,
$p = 5.9 \times 10^{-7}$). Predicted slope for constant lifetime
burden: $-1.0$ (blue dotted line).
\textbf{B.}~Lifetime mutation burden (rate $\times$ lifespan) is
approximately constant across species (mean $= 2{,}191$,
CV $= 0.34$). Red band: $\pm 1$ SD.
\textbf{C.}~Confidence per governance layer (orders of suppression /
layer count) vs.\ lifespan for the three species with both Cagan
somatic rates and governance layer data. Red ring intensity indicates
cancer rate.}
\label{fig:cagan}
\end{figure}

\subsection{The constant}

If governance confidence is calibrated to lifespan, then:
\[
\mu_{\text{somatic}} \times L \approx T
\]
where $\mu_{\text{somatic}}$ is the somatic mutation rate per year,
$L$ is lifespan, and $T$ is the tolerable lifetime burden. Across
16 mammals, $T \approx 2{,}191$ (mean; CV $= 0.34$). The slope of
$\log(\mu)$ vs.\ $\log(L)$ is $-0.81$, close to the $-1.0$ predicted
by a perfect constant. The residual (slope $\neq -1$) indicates that
longer-lived species tolerate slightly more total burden, likely due
to greater redundancy in downstream repair.

This means governance confidence directly predicts lifespan:
\[
L \approx \frac{T}{\mu_{\text{somatic}}} \propto \text{confidence}
\]

Not through a proxy. Through the system output of the governance
product. Mouse: 796 SBS/year, 2.7-year lifespan, 0.32 orders of
suppression per layer. Naked mole rat: 93 SBS/year, 37-year lifespan,
0.43 orders per layer. Human: 47 SBS/year, 80-year lifespan, 0.47
orders per layer. The ranking is exact.

\subsection{Two optimization surfaces}

The Cagan data covers mammals. Mammals are determinate growers: cell
count is fixed at maturity. The constant $T$ is mutations per cell at
end of life.

Reptiles, fish, and amphibians are often indeterminate growers:
cell count increases throughout life. They do not stop growing.
They maintain continuous telomerase expression in somatic cells
(mammals shut it off). Despite this, their cancer rates are lower:
mammals $\sim$12\% at death, reptiles and birds $\sim$4\%,
amphibians $\sim$1.2\% \cite{Vincze2024}.

The fixed metric for indeterminate growers cannot be total burden. It
must be mutation \emph{density}: burden per cell at any given time.
The governance is maintaining a rate sustainable as the system scales.

\begin{center}
\small
\begin{tabular}{lll}
\toprule
& \textbf{Mammals} & \textbf{Indeterminate growers} \\
\midrule
Optimization target & Context per cell & System longevity \\
Cell count & Fixed at maturity & Growing \\
Fixed metric & Burden/cell at death & Burden density/cell/time \\
Telomerase & Off (limits divisions) & On (continuous renewal) \\
Specialization & Extreme & Moderate \\
Cancer strategy & Hayflick limit & Maintain governance at scale \\
\bottomrule
\end{tabular}
\end{center}

Mammals traded system longevity for context per cell. Reptiles traded
context per cell for system longevity. A crocodile neuron is simpler
than a human neuron. A tortoise liver cell is less differentiated than
a human hepatocyte. The governance overhead per cell is lower, which
is why the system can keep adding cells indefinitely.

The Greenland shark lives 400+ years with 81 duplicated DNA repair
genes and an altered TP53 \cite{Spath2024}. Galapagos tortoises have
duplicated tumor suppressors (NF2, SMAD4, PML) and 75\% of
turtle species show negligible senescence \cite{daSilva2022}. These
organisms did not solve the governance problem by increasing
confidence per layer (the mammalian solution). They solved it by
keeping the governance overhead low enough that continuous growth
remains viable.

Cephalopods took a third path entirely: extensive RNA editing
(11,000+ active sites \cite{Liscovitch2017}) instead of DNA repair for adaptive flexibility.
They sacrifice genomic evolution rate to preserve editing sites. Short
lifespan (1--5 years), but extraordinary per-cell computational
flexibility. All context, no longevity. The inverse of the reptile
strategy.

No somatic SBS/year data exists for any non-mammalian vertebrate.
This is the critical measurement gap. The framework predicts that
each vertebrate class has its own constant $T$, determined by
its governance architecture. Within each class, the Cagan inverse
should hold. Across classes, the constants should differ.

% ------------------------------------------------------------------
\section{The Monolith Tradeoff}
\label{sec:monolith}
% ------------------------------------------------------------------

A monolith with 2 layers at 99.999\% each delivers higher availability
than a microservice mesh with 7 layers at 99.9\%. The product math
favors fewer layers \emph{if} each layer is near-perfect.

Bacteria are the monolith. Two repair layers. Ancient. Refined over
3.8 billion years. Extremely high confidence per layer. No coordination
overhead. No service mesh. And it works: bacteria do not get cancer.

But a monolith cannot scale vertically. The same reason a single-process
application handles 100 requests per second and collapses at 100,000.
Bacteria solved this by horizontal scaling: stay small, replicate fast,
let the colony absorb individual failures.

The moment the organism needs $3.7 \times 10^{13}$ cells to
coordinate, it needs the mesh. And the mesh introduces:

\begin{itemize}
\item Each layer is a new attack surface (MMR genes can themselves mutate)
\item Layers interact (DSB repair depends on checkpoints for time)
\item Coordination has overhead (immune surveillance costs energy,
   causes autoimmunity)
\item The mesh needs an orchestrator (p53), which becomes a single
   point of failure
\end{itemize}

TP53 is mutated in over half of all human cancers not because it is
fragile but because it is the orchestrator. The most attacked node in
any distributed system is the load balancer.

% ------------------------------------------------------------------
\section{Predictions}
\label{sec:predictions}
% ------------------------------------------------------------------

The framework generates testable predictions:

\begin{enumerate}
\item \textbf{Layer addition threshold.} If governance layers scale as
   $\log(n)$, there should be a cell count threshold at which a new
   layer becomes necessary. Organisms near $10^{16}$ cells (blue whales)
   should show evidence of an eighth layer or enhanced redundancy in
   existing layers.

\item \textbf{Cancer as layer readout.} The mutation rate in a tumor
   should predict which governance layers remain functional. MSI-high
   tumors ($\sim 10^{-7}$) have lost MMR. Tumors at $\sim 10^{-5}$
   have lost multiple layers. The rate is a count of remaining layers.

\item \textbf{CHIP precedes cancer by the WAL-to-data delay.} CHIP
   mutations (WAL degradation) should precede driver mutations (data
   corruption) by the number of cell divisions required for index
   loss to produce data loss. This lag is observed: CHIP is detectable
   years to decades before cancer diagnosis \cite{Jaiswal2014}.

\item \textbf{Layer confidence predicts cancer spectrum.} Organisms
   with the same layer count but different per-layer confidence should
   show different cancer spectra. The naked mole rat (enhanced
   checkpoints) should be resistant to checkpoint-driven cancers
   specifically. Evidence: the rare naked mole rat tumors that have
   been documented are \emph{not} carcinomas (checkpoint-driven) but
   sarcomas (mesenchymal, potentially driven by other layers)
   \cite{Delaney2016}.

\item \textbf{The weakest layer predicts the first cancer.} In a
   given species, the most common cancer type should correspond to
   the lowest-confidence governance layer. In mice (low checkpoint
   confidence), lymphoma dominates. In humans (low WAL confidence),
   the cancers that increase most with age are those in rapidly
   dividing tissues (colon, skin, blood) where WAL burden is highest.

\item \textbf{Cross-class constants.} The Cagan inverse relationship
   ($\mu \times L \approx T$) should hold within each vertebrate
   class but with a different constant $T$ per class. Mammals
   ($T \approx 2{,}200$), reptiles (predicted: lower $T$, reflecting
   lower per-cell governance burden), and birds (predicted: higher $T$,
   reflecting their elevated metabolic rate and enhanced DSB repair).
   No somatic SBS/year data exists for non-mammals. This measurement
   would directly test the framework.

\item \textbf{Indeterminate growers optimize density, not total.}
   For reptiles and fish with continuous growth, the fixed metric
   should be mutation burden per cell (density), not total mutations.
   Plotting mutations/cell vs.\ age should yield a constant slope
   regardless of body size, while total mutations should scale with
   size. For mammals, the opposite: total burden is constant, but
   burden per cell is irrelevant because cell count is fixed.

\item \textbf{Early-diverging mammals as intermediates.} Marsupials
   and monotremes diverged from placental mammals $\sim$160~MYA.
   They should show intermediate governance profiles: mammalian layer
   count but with different per-layer confidence, reflecting their
   older governance architecture. Opossums show $\sim$56\% neoplasia
   rate \cite{Vincze2024}, the highest of any mammal studied. The
   governance version is outdated.
\end{enumerate}

% ------------------------------------------------------------------
\section{Relationship to the Series}
\label{sec:series}
% ------------------------------------------------------------------

This paper is the existence proof. Papers~1--4 derive from first
principles that a persistent system under bounded context must
maintain a directed graph, must escalate, must couple knowledge with
action, and must encode at permanence levels proportional to
confidence. Papers~5--8 show that cancer data is consistent with
these constraints.

This paper shows the genome doing all of it.

\begin{center}
\small
\begin{tabular}{lll}
\toprule
\textbf{Paper} & \textbf{Prediction} & \textbf{Genome instantiation} \\
\midrule
1 (Graph) & Must maintain directed graph & DNA + regulatory network \\
2 (Escalation) & Top-down allocation via escalation & 7-layer repair hierarchy \\
3 (K/A) & Knowledge and action inseparable & Coding + regulatory in same sequence \\
4 (Encoding) & Permanence $\propto$ confidence & DNA $>$ epigenetics $>$ protein state \\
\midrule
5 (Cancer) & Disruption follows graph topology & Mutations cluster by channel \\
6 (Channels) & Channels carry independent signal & Channel count predicts survival \\
7 (Position) & Connectivity predicts severity & Degree $\times$ architecture interaction \\
8 (Context) & Context diseases share machinery & Chromatin genes serve multiple functions \\
\bottomrule
\end{tabular}
\end{center}

The genome did not read these papers. It arrived at the same
architecture because it is subject to the same constraints.
3.8 billion years of selection under bounded context produced the
same design patterns that software engineering rediscovered in 60
years.

% ------------------------------------------------------------------
\section{Discussion}
\label{sec:discussion}
% ------------------------------------------------------------------

\subsection{What this paper does not claim}

This paper does not claim the genome was ``designed'' as a knowledge
graph. It claims that selection under bounded context produces graph
structure as a necessary consequence (Paper~1), and that the genome
is an instance of this necessity.

This paper does not claim cancer is ``merely'' an engineering problem.
It claims that the structural constraints governing cancer are the
same structural constraints governing any distributed system under
bounded context. The analogy is not metaphorical. It is mathematical.

\subsection{The measurement gap}

The central limitation is per-layer confidence measurement. We can
count layers. We can measure system output. We cannot decompose
confidence per layer without either destructive testing (knockouts)
or indirect inference (cancer spectrum analysis). Closing this gap
would enable the $n$-nines decomposition that distributed systems
engineering takes for granted.

Cancer genomics is beginning to close it. Every sequenced tumor is a
natural knockout experiment. Systematic comparison of mutation
signatures across tumor types, mapped to known repair pathway
deficiencies, could yield per-layer confidence estimates across
species and tissue types.

\subsection{Implications for cancer treatment}

If cancer is an autoscaling failure, treatment should restore the
autoscaling policy, not merely kill the rogue instances. Chemotherapy
kills rogue instances. Immunotherapy restores the monitoring layer.
Targeted therapy (PARP inhibitors, checkpoint inhibitors) restores
specific governance layers. The framework predicts that treatments
targeting the \emph{specific failed layer} should outperform treatments
that target the tumor generically. Early evidence supports this:
PARP inhibitors in BRCA-mutant cancers (restoring DSB repair
governance) show higher efficacy than in BRCA-wild-type cancers
\cite{Lord2017}.

% ------------------------------------------------------------------
\section{Conclusion}
\label{sec:conclusion}
% ------------------------------------------------------------------

The genome is a knowledge graph under bounded context. DNA is its
serialized projection. Epigenetic marks are its write-ahead log.
Repair layers are its governance architecture. Their number scales
logarithmically with the system they govern.

Cancer is not a disease of uncontrolled growth. It is a disease of
uncontrolled growth \emph{control}. The cell's capacity to replicate
is not the pathology---it is the substrate on which the organism
depends. The pathology is in the policy, not the process.

In a bounded system, every resource consumed by uncoordinated function
is unavailable to coordinated function. The tumor does not need to be
malicious. It only needs to be uncoordinated. That is sufficient.

\section*{Acknowledgments}

The author used Claude (Anthropic) for drafting assistance, data
compilation, and \LaTeX{} preparation. All theoretical content and
arguments are the author's own.

\bibliographystyle{plain}
\begin{thebibliography}{99}

\bibitem{McCarthy2026a}
P.~D. McCarthy.
\newblock Graph structure is necessary for information preservation
under bounded context.
\newblock \emph{SSRN preprint}, 2026.

\bibitem{McCarthy2026d}
P.~D. McCarthy.
\newblock Evaluation-driven state revision and the encoding hierarchy
under bounded context.
\newblock \emph{SSRN preprint}, 2026.

\bibitem{Jaiswal2014}
S.~Jaiswal, P.~Fontanillas, J.~Flannick, et~al.
\newblock Age-related clonal hematopoiesis associated with adverse
outcomes.
\newblock \emph{New England Journal of Medicine}, 371(26):2488--2498, 2014.

\bibitem{Seluanov2009}
A.~Seluanov, C.~Hine, M.~Bozzella, et~al.
\newblock Distinct tumor suppressor mechanisms evolve in rodent species
that differ in size and lifespan.
\newblock \emph{Aging Cell}, 8(4):353--361, 2009.

\bibitem{Sulak2016}
M.~Sulak, L.~Fong, K.~Mika, et~al.
\newblock TP53 copy number expansion is associated with the evolution
of increased body size and an enhanced DNA damage response in elephants.
\newblock \emph{eLife}, 5:e11994, 2016.

\bibitem{Hanahan2011}
D.~Hanahan and R.~A. Weinberg.
\newblock Hallmarks of cancer: the next generation.
\newblock \emph{Cell}, 144(5):646--674, 2011.

\bibitem{Delaney2016}
M.~A. Delaney, L.~Ward, T.~Walsh, et~al.
\newblock Initial case reports of cancer in naked mole-rats
(\textit{Heterocephalus glaber}).
\newblock \emph{Veterinary Pathology}, 53(3):691--696, 2016.

\bibitem{Lord2017}
C.~J. Lord and A.~Ashworth.
\newblock PARP inhibitors: synthetic lethality in the clinic.
\newblock \emph{Science}, 355(6330):1152--1158, 2017.

\bibitem{Cagan2022}
A.~Cagan, A.~Baez-Ortega, N.~Brzozowska, et~al.
\newblock Somatic mutation rates scale with lifespan across mammals.
\newblock \emph{Nature}, 604(7906):517--524, 2022.

\bibitem{Drake1991}
J.~W. Drake.
\newblock A constant rate of spontaneous mutation in DNA-based microbes.
\newblock \emph{Proceedings of the National Academy of Sciences},
88(16):7160--7164, 1991.

\bibitem{Lynch2010}
M.~Lynch.
\newblock Evolution of the mutation rate.
\newblock \emph{Trends in Genetics}, 26(8):345--352, 2010.

\bibitem{Bianconi2013}
E.~Bianconi, A.~Piovesan, F.~Facchin, et~al.
\newblock An estimation of the number of cells in the human body.
\newblock \emph{Annals of Human Biology}, 40(6):463--471, 2013.

\bibitem{Vincze2024}
O.~Vincze, F.~Colchero, J.-F. Lema\^itre, et~al.
\newblock Cancer prevalence across vertebrates.
\newblock \emph{Cancer Discovery}, 15(1):227--237, 2024.

\bibitem{Spath2024}
P.~Spath, E.~Kuderna, A.~Valen, et~al.
\newblock The Greenland shark genome reveals longevity-associated
DNA repair gene expansions.
\newblock \emph{bioRxiv}, 2024.

\bibitem{daSilva2022}
R.~da~Silva, D.~A. Conde, A.~Baudisch, and F.~Colchero.
\newblock Slow and negligible senescence among testudines challenges
evolutionary theories of senescence.
\newblock \emph{Science}, 376(6600):1466--1470, 2022.

\bibitem{Liscovitch2017}
N.~Liscovitch-Brauer, S.~Alon, H.~T. Porath, et~al.
\newblock Trade-off between transcriptome plasticity and genome evolution in
cephalopods.
\newblock \emph{Cell}, 169(2):191--202, 2017.

\bibitem{VanderHeiden2009}
M.~G. {Vander Heiden}, L.~C. Cantley, and C.~B. Thompson.
\newblock Understanding the Warburg effect: the metabolic requirements of cell
proliferation.
\newblock \emph{Science}, 324(5930):1029--1033, 2009.

\end{thebibliography}

\end{document}
